\subsection{工作流引导、科学判断与展望}

50题对照中，Guided的平均自洽性和准确度均高于Standard；Hard中，两者对增益的加权贡献分别为+0.1353和+0.1274。HH1和HI1展示了IRC及端点计算如何确定接近参考能垒的鞍点实际连接哪些物种。

HA8和NU6展示了如何依据计算结构与路径确定能量参照和鞍点归属。HA8采用稳定反应物极小值作为能垒参照，而IRC连接的是鞍点附近的局部活性构象。NU6中，用于寻找betaine至产物闭环鞍点的初猜反复优化至直接反应物至产物的鞍点区域。DAG同时保存反应物构象、鞍点、IRC端点及相应步骤，使智能体能在新计算返回后检查每个能量属于哪个结构。计算结果由此成为智能体规划和检验后续研究的外部反馈。\citep{valmeekam2023planning}

这些情况需要不同的后续操作：补算缺少的端点，调整反复回到另一条路线的初猜，或重新评估无法描述目标电子态的方法。计算结构或路径改变后，智能体应相应更新其反应描述。\citep{turpin2023unfaithful} Record保留了工具，自洽分仍低于Guided，Normal和Hard任务尤其如此。归档轨迹展示了工具反馈如何引导下一次几何操作或计算。\citep{ref18,huang2024selfcorrect}

Brief-errors在IA1和IA9中得分更高，Guided则在IG8中更高。这些交叉结果提出了一个依赖任务的假说：详细诊断可能有助于修复TS或IRC计算；方法或目标本身需要重新审视时，继续局部修补的收益可能有限。

G2中，反复检查促使智能体删除了原生计算已支持的数值。Supervisor反馈曾恢复这些字段，后续拒绝又导致最终删除。检验路径归属时，应在机理图中保留数值观测及其日志。

P2每题每配置保留一次运行，任务簇bootstrap区间反映题目之间的变化。P1b在四题、四配置的每个组合上重复三次，共48次运行。部分Hard任务的Accuracy目标是定性通道身份，因此得分提高同时反映路径识别、验证计算的补全和能垒数值的差异。

Standard有时在计算仍进行时结束研究，Guided/HI1则在收到完成通知后仍停留于等待状态。计算结果返回后，应先恢复尚未完成的研究，再形成最终报告。NU6反复收敛至同一鞍点区域的情况，也提示应比较新候选与先前尝试，再选择下一次初猜。

默认预组织采样32个刚体相对放置，进行局部MMFF94松弛，再按能量与聚类选择候选。CREST通过可选后端提供额外采样。

本基准以电子能差计分，包括弱酸性水溶液任务NI5和高浓度二氯甲烷任务NU5所规定的电子能参照。溶液条件下的$\Delta G^\ddagger$还涉及零点能、热与熵、标准态及溶剂化贡献；溶剂参与也可能改变描述路径所需的化学模型。溶剂敏感性仍待评估。

各工作流采用固定的原子映射与分支拥挤设置，并使用DeepSeek-v4.1-flash（1017），API别名为\texttt{deepseek-v4-flash}。

Cloudra以GNU Affero General Public License（AGPL）开源，代码公开于\url{https://github.com/yyywaa/Cloudra-mechanism-agent}，支持系统的公开检查、复用与社区共同发展。可复现的表格与分数见配套数据包。
